\chapter{Introducción}

El mercado financiero, entendido como el espacio en el cual se compran y se venden activos financieros (acciones, bonos, divisas, materias primas y criptomonedas principalmente), constituye un ámbito de estudio especialmente adecuado para poner en práctica el uso de tecnologías Big Data.

Las tecnologías Big Data se definen como las herramientas que permiten gestionar, almacenar y analizar conjuntos de datos masivos que superan la capacidad de tratamiento tradicional.

Desde los comienzos de la teoría económica moderna, se ha considerado el mercado financiero como un sistema caótico. A diferencia de un sistema puramente aleatorio, un sistema caótico presenta patrones y reglas subyacentes que permiten realizar determinadas predicciones sobre su comportamiento.

El caos se caracteriza por la complejidad. En un sistema caótico, las variables presentan relaciones no lineales, a menudo difíciles de detectar a simple vista. Esto deriva en la impredecibilidad del sistema, puesto que pequeñas variaciones en las condiciones iniciales pueden llevar a resultados completamente diferentes.

Además, el mercado financiero puede entenderse como un sistema de segundo orden, dado que las propias predicciones que se realizan sobre él pueden influir en los resultados que se obtienen.

En el presente trabajo se propone abordar el problema de confeccionar una cartera de inversión, esto es, un subconjunto de los activos financieros disponibles, de forma eficiente.

\cite{markowitz1952} introduce el modelo de media-varianza como una de las aproximaciones clásicas para la gestión de carteras de inversión. Este modelo se fundamenta en el principio riesgo-rendimiento, buscando maximizar el rendimiento esperado para un determinado nivel de riesgo asumido por el inversor.

Con el uso de tecnologías Big Data y, en particular, con el uso de modelos generativos como los \textbf{Autoencoders Variacionales (VAEs)} y las \textbf{Redes Generativas Adversarias (GANs)}, se propone una aproximación a la confección de carteras desde una perspectiva no lineal.

\section{Motivación}

La gestión de carteras de inversión ha estado dominada durante décadas por la Teoría de Cartera Moderna (MPT), introducida por \cite{markowitz1952}. Este marco teórico se fundamenta en la optimización media-varianza, bajo el supuesto de que los retornos de los activos financieros siguen una distribución normal. Sin embargo, la evidencia empírica en el análisis de datos masivos financieros demuestra que los mercados reales presentan características que este modelo no logra capturar: la curtosis excesiva (colas pesadas), la asimetría y la volatilidad estocástica.

En el contexto actual de Big Data, la incapacidad de los modelos lineales para gestionar la complejidad representa una limitación relevante. Los eventos extremos ocurren con una frecuencia mayor en la realidad de la que predeciría un modelo normal, lo que hace que la varianza sea una métrica de riesgo insuficiente. El uso de modelos generativos como las GANs \cite{NIPS2016_8a3363ab} y los VAEs \cite{kingma2013auto} ofrece un nuevo paradigma para aprender la distribución subyacente de los datos sin imponer una forma paramétrica rígida.

Frente a otras aproximaciones, como el uso de modelos de Machine Learning o sistemas de ecuaciones diferenciales, los modelos generativos permiten trabajar desde una perspectiva probabilística. Estos modelos se adaptan a los datos y ofrecen flexibilidad frente a cambios drásticos.

\section{Planteamiento del trabajo}

La elección de esta arquitectura, como se exponía previamente, no busca predecir el precio futuro de los activos, sino generar diversos escenarios posibles para comprender el riesgo subyacente de los mismos.

El presente Trabajo de Fin de Máster propone una metodología híbrida que integra el modelado generativo y la teoría de la información. El núcleo de la investigación se divide en cuatro fases:

\begin{enumerate}
\item Extracción de datos de series históricas del S\&P 500. Se propone la creación de un pipeline de datos que permita la implementación de la arquitectura propuesta.
\item Uso de \textbf{Autoencoders Variacionales (VAEs)} para la reducción de dimensionalidad no lineal.
\item Desarrollo de una \textbf{Red Generativa Adversaria (GAN)} para sintetizar escenarios de mercado realistas.
\item Optimización de la cartera mediante la \textbf{Divergencia de Kullback-Leibler} como métrica de riesgo informativo.
\end{enumerate}

\section{Estructura del trabajo}

La memoria se organiza de la siguiente manera: el Capítulo 2 presenta los objetivos del trabajo, diferenciando entre el objetivo general, los objetivos específicos y el alcance de la investigación. El Capítulo 3 desarrolla el marco teórico, introduciendo la teoría moderna de carteras, el modelo de media-varianza de Markowitz, la teoría de la información y los principales modelos generativos considerados. El Capítulo 4 revisa el estado del arte, analizando los enfoques existentes en optimización de carteras, aprendizaje automático aplicado a finanzas y modelado generativo de series financieras. El Capítulo 5 describe la metodología propuesta, incluyendo la adquisición y el preprocesamiento de datos, la arquitectura VAE-GAN, la generación de escenarios sintéticos y los criterios de comparación. El Capítulo 6 recoge las herramientas y tecnologías empleadas. El Capítulo 7 expone el marco normativo, ético y de protección de datos. El Capítulo 8 presenta la planificación del trabajo. Finalmente, el Capítulo 9 recoge las conclusiones del estudio.
